% resultater.tex --------------------------------

% Kort kapittelintroduksjon

\subsection{FFT-resultater}
% Presentasjon av resultatene fra Fourier-transformasjonen (hele signalet)
% Beskriv analysens forutsetninger: varighet, samplingsrate og eventuelle filtre

\subsubsection{Spektrumanalyse}
% Presenter resultatet fra FFT-beregningen som figur eller tabell
% Vis det observerte frekvensspekteret og marker de mest fremtredende toppene
% Kommenter kun hvilke frekvenser som forekommer og deres relative styrke
% Unngå forklaring av årsaker – det skal gjøres i diskusjonen

\subsubsection{Harmoniske komponenter}
% Identifiser grunnfrekvens og overtoner i spekteret
% Angi målte verdier i Hz og sammenlign med teoretiske frekvenser uten å vurdere avvik
% Eventuelt bruk tabell for å liste grunnfrekvens + 2–3 harmoniske overtoner

\subsubsection{Amplitude og energifordeling}
% Presenter hvordan signalenergien er fordelt mellom ulike frekvensbånd
% Bruk graf eller tabell for å vise total energi i lav-, mellom- og høyfrekvent område
% Beskriv observasjonene nøytralt uten å forklare årsaker




\subsection{STFT-resultater}
% Resultater fra tids-frekvensanalyse (spektrogram)
% Forklar kort hvilke parametere som ble brukt (vindutype, størrelse, overlapp)

\subsubsection{Spektrogram}
% Vis spektrogrammet som figur med akser for tid og frekvens
% Beskriv de tydelige mønstrene og hvor i tidsaksen de opptrer
% Eksempel: «Tydelige bånd observeres mellom 50–110 Hz gjennom hele klippet»

\subsubsection{Tidsvariasjon og rytmikk}
% Kommenter hvordan spektrogrammet viser variasjoner i tid og rytmiske elementer
% Beskriv observerte mønstre, f.eks. endringer i amplitude eller hyppighet av slag
% Ikke forklar hvorfor variasjonene oppstår – det tas opp i diskusjon.tex





\subsection{Oppsummering}
% Kort oppsummering av hovedfunn fra FFT- og STFT-analysen
% Nevn hvilke frekvensområder og mønstre som ble observert
% Avslutt med en setning som peker frem mot diskusjonsdelen, f.eks.:
% «I neste kapittel drøftes hvordan disse resultatene kan forklares ut fra teori og metodevalg.»
